\documentclass{beamer}

\usepackage{amsmath,amsfonts,amssymb}
\usepackage{bm}

\title{Quantum Approximate Optimization Algorithm (QAOA)}
\author{Morten Hjorth-Jensen}
\date{Spring 2026}

\begin{document}

\frame{\titlepage}

%------------------------------------------------
\section{Background}
%------------------------------------------------

\begin{frame}{Combinatorial Optimization Problems}
\begin{itemize}
\item Goal: Optimize an objective function over discrete variables
\[
C(x): \{0,1\}^n \rightarrow \mathbb{R}
\]
\item Many problems (MaxCut, Knapsack, TSP) are NP-hard
\item Exact solutions scale exponentially with system size
\item Approximate algorithms aim for
\[
\alpha = \frac{C(x^*)}{C_{\max}} \lesssim 1
\]
\end{itemize}
\end{frame}

%------------------------------------------------

\begin{frame}{QUBO Formulation}
\begin{itemize}
\item Quadratic Unconstrained Binary Optimization (QUBO):
\[
C(x) = x^T Q x = \sum_{i,j} Q_{ij} x_i x_j
\]
\item Optimization problem:
\[
x^* = \arg\min_{x \in \{0,1\}^n} C(x)
\]
\item NP-hard in general
\item Can be mapped to Ising variables:
\[
z_i = 2x_i - 1, \quad z_i \in \{-1,1\}
\]
\end{itemize}
\end{frame}

%------------------------------------------------

\begin{frame}{QUBO and Ising Hamiltonians}
\begin{itemize}
\item QUBO $\leftrightarrow$ Ising model mapping
\item Cost function becomes Hamiltonian:
\[
\hat{H}_C = \sum_{i,j} J_{ij} Z_i Z_j + \sum_i h_i Z_i
\]
\item Enables quantum optimization via ground state search
\item Basis for both annealing and QAOA
\end{itemize}
\end{frame}

%------------------------------------------------

\begin{frame}{MaxCut Problem}
\begin{itemize}
\item Graph $G=(V,E)$ with weights $w_{ij}$
\item Partition vertices into two sets
\item Objective:
\[
C(x) = \sum_{i,j} w_{ij} x_i(1-x_j)
\]
\item NP-hard problem
\item Benchmark problem for QAOA
\end{itemize}
\end{frame}

%------------------------------------------------

\begin{frame}{Classical Optimization Methods}
\begin{itemize}
\item Exact methods:
\begin{itemize}
\item Branch-and-bound
\item Quadratic programming
\end{itemize}
\item Approximation methods:
\begin{itemize}
\item Semidefinite programming (Goemans–Williamson)
\item Approximation ratio $\alpha \approx 0.878$
\end{itemize}
\item Heuristics:
\begin{itemize}
\item Simulated annealing
\item Genetic algorithms
\end{itemize}
\end{itemize}
\end{frame}

%------------------------------------------------

\begin{frame}{Variational Quantum Algorithms (VQAs)}
\begin{itemize}
\item Hybrid quantum-classical optimization loop
\item Parameterized quantum state:
\[
|\psi(\theta)\rangle = U(\theta)|0\rangle
\]
\item Variational principle:
\[
E(\theta) = \langle \psi(\theta) | \hat{H} | \psi(\theta) \rangle
\]
\item Classical optimizer minimizes $E(\theta)$
\end{itemize}
\end{frame}

%------------------------------------------------

\begin{frame}{Variational Quantum Eigensolver (VQE)}
\begin{itemize}
\item Solve eigenvalue problem:
\[
E_0 \approx \min_{\theta} \langle \psi(\theta)|\hat{H}|\psi(\theta)\rangle
\]
\item Key ingredients:
\begin{itemize}
\item Ansatz $U(\theta)$
\item Measurement of expectation values
\item Classical optimization
\end{itemize}
\end{itemize}
\end{frame}

%------------------------------------------------

\begin{frame}{Quantum Adiabatic Algorithm (QAA)}
\begin{itemize}
\item Time-dependent Hamiltonian:
\[
\hat{H}(t) = f(t)\hat{H}_C + g(t)\hat{H}_M
\]
\item Adiabatic theorem:
\begin{itemize}
\item Slow evolution preserves eigenstate
\end{itemize}
\item Trotterized evolution:
\[
U(t) \approx \prod_k e^{-i \hat{H}_C \Delta t} e^{-i \hat{H}_M \Delta t}
\]
\end{itemize}
\end{frame}

%------------------------------------------------

\begin{frame}{QAOA Overview}
\begin{itemize}
\item Discretized version of QAA
\item Alternating layers:
\[
|\psi_p(\gamma,\beta)\rangle =
\prod_{k=1}^p e^{-i\beta_k \hat{H}_M} e^{-i\gamma_k \hat{H}_C} |s\rangle
\]
\item Initial state:
\[
|s\rangle = |+\rangle^{\otimes n}
\]
\item Optimize parameters $(\gamma,\beta)$
\end{itemize}
\end{frame}

%------------------------------------------------

\begin{frame}{QAOA Cost Function}
\begin{itemize}
\item Expectation value:
\[
F_p(\gamma,\beta) =
\langle \psi_p | \hat{H}_C | \psi_p \rangle
\]
\item Optimization:
\[
(\gamma^*,\beta^*) = \arg\max F_p
\]
\item Approximation ratio:
\[
\alpha = \frac{F_p(\gamma^*,\beta^*)}{C_{\max}}
\]
\end{itemize}
\end{frame}

%------------------------------------------------
\section{QAOA Analysis}
%------------------------------------------------

\begin{frame}{QAOA Ansatz Structure}
\begin{itemize}
\item General unitary:
\[
U(\gamma,\beta) =
\prod_{k=1}^p e^{-i\beta_k \hat{H}_M} e^{-i\gamma_k \hat{H}_C}
\]
\item Choice of ansatz affects:
\begin{itemize}
\item Expressivity
\item Trainability
\item Hardware efficiency
\end{itemize}
\end{itemize}
\end{frame}

%------------------------------------------------

\begin{frame}{Ansatz Variants}
\begin{itemize}
\item Multi-angle QAOA:
\begin{itemize}
\item Independent parameters per gate
\end{itemize}
\item Adaptive QAOA:
\begin{itemize}
\item Problem-specific ansatz construction
\end{itemize}
\item Counterdiabatic QAOA:
\begin{itemize}
\item Faster convergence via additional terms
\end{itemize}
\end{itemize}
\end{frame}

%------------------------------------------------

\begin{frame}{Parameter Optimization}
\begin{itemize}
\item Classical optimization challenges:
\begin{itemize}
\item High-dimensional parameter space
\item Local minima
\end{itemize}
\item Strategies:
\begin{itemize}
\item Good initialization
\item Gradient-based methods
\item Exploiting symmetries
\end{itemize}
\end{itemize}
\end{frame}

%------------------------------------------------

\begin{frame}{Barren Plateaus}
\begin{itemize}
\item Problem: exponentially vanishing gradients
\[
\text{Var}[\partial_i C(\theta)] \sim \mathcal{O}(b^{-n})
\]
\item Consequence:
\begin{itemize}
\item Training becomes inefficient
\end{itemize}
\item Mitigation:
\begin{itemize}
\item Structured ansätze
\item Problem-inspired circuits
\end{itemize}
\end{itemize}
\end{frame}

%------------------------------------------------

\begin{frame}{Computational Efficiency}
\begin{itemize}
\item Potential quantum speedup unclear
\item Limitations:
\begin{itemize}
\item Circuit depth
\item Optimization overhead
\end{itemize}
\item Improvements:
\begin{itemize}
\item Better ansätze
\item Parameter reuse
\end{itemize}
\end{itemize}
\end{frame}

%------------------------------------------------

\begin{frame}{Solution Quality}
\begin{itemize}
\item Performance measured by approximation ratio
\item Depends on:
\begin{itemize}
\item Circuit depth $p$
\item Problem structure
\end{itemize}
\item $p \to \infty$:
\begin{itemize}
\item Recovers adiabatic evolution
\end{itemize}
\end{itemize}
\end{frame}

%------------------------------------------------

\begin{frame}{Noise and Hardware Effects}
\begin{itemize}
\item NISQ devices introduce:
\begin{itemize}
\item Gate errors
\item Decoherence
\end{itemize}
\item Strategies:
\begin{itemize}
\item Noise mitigation
\item Hardware-aware ansätze
\end{itemize}
\end{itemize}
\end{frame}

%------------------------------------------------

\begin{frame}{Summary}
\begin{itemize}
\item QAOA is a leading variational algorithm for optimization
\item Built on:
\begin{itemize}
\item QUBO/Ising mapping
\item Hybrid optimization
\end{itemize}
\item Key challenges:
\begin{itemize}
\item Parameter optimization
\item Noise and scalability
\end{itemize}
\item Active research area with many variants
\end{itemize}
\end{frame}

%------------------------------------------------

\end{document}



%================================================
\section{QAOA for MaxCut: Explicit Construction}
%================================================

\begin{frame}{MaxCut as a Quantum Hamiltonian}
\begin{itemize}
\item Classical objective:
\[
C(x) = \sum_{(i,j)\in E} w_{ij} x_i (1 - x_j)
\]
\item Map binary variables:
\[
x_i = \frac{1 - Z_i}{2}
\]
\item Cost Hamiltonian:
\[
\hat{H}_C = \frac{1}{2} \sum_{(i,j)\in E} w_{ij}(I - Z_i Z_j)
\]
\item Ground state encodes optimal cut
\end{itemize}
\end{frame}

%------------------------------------------------

\begin{frame}{Mixer Hamiltonian}
\begin{itemize}
\item Standard choice:
\[
\hat{H}_M = \sum_i X_i
\]
\item Properties:
\begin{itemize}
\item Does not commute with $\hat{H}_C$
\item Generates transitions between bitstrings
\end{itemize}
\item Initial state:
\[
|s\rangle = |+\rangle^{\otimes n}
\]
\end{itemize}
\end{frame}

%------------------------------------------------

\begin{frame}{QAOA Circuit Structure}
\begin{itemize}
\item One QAOA layer:
\[
U_C(\gamma) = e^{-i\gamma \hat{H}_C}, \quad
U_M(\beta) = e^{-i\beta \hat{H}_M}
\]
\item State:
\[
|\psi_p\rangle = \prod_{k=1}^p U_M(\beta_k) U_C(\gamma_k) |s\rangle
\]
\end{itemize}
\end{frame}



\begin{frame}{Gate Decomposition (MaxCut)}
  \begin{itemize}
  \item Two-qubit interaction:
    \[
    e^{-i\gamma Z_i Z_j}
    \]
  \item Decomposition:
    \begin{itemize}
    \item CNOT -- $R_Z$ -- CNOT
    \end{itemize}
  \item Mixer:
    \[
    e^{-i\beta X_i} = R_X(2\beta)
    \]
  \end{itemize}
\end{frame}

%================================================
\section{Analytic Solution for $p=1$}
%================================================

\begin{frame}{QAOA at $p=1$}
  \begin{itemize}
  \item State:
    \[
    |\psi(\gamma,\beta)\rangle = e^{-i\beta \hat{H}_M} e^{-i\gamma \hat{H}_C} |+\rangle^{\otimes n}
    \]
  \item Expectation value:
    \[
    F(\gamma,\beta) =
    \langle \psi | \hat{H}_C | \psi \rangle
    \]
  \end{itemize}
\end{frame}

%------------------------------------------------

\begin{frame}{Edge Contribution}
  \begin{itemize}
  \item For one edge $(i,j)$:
    \[
    \langle Z_i Z_j \rangle =
    \sin(4\beta)\sin(2\gamma)
    \]
  \item Total energy:
    \[
    F(\gamma,\beta) =
    \sum_{(i,j)} \frac{w_{ij}}{2}
    \left(1 - \sin(4\beta)\sin(2\gamma)\right)
    \]
  \end{itemize}
\end{frame}

%------------------------------------------------

\begin{frame}{Optimal Parameters (Regular Graphs)}
  \begin{itemize}
  \item For regular graphs:
    \[
    \beta^* = \frac{\pi}{8}, \quad
    \gamma^* = \frac{\pi}{4}
    \]
  \item Leads to approximation ratio:
    \[
    \alpha \approx 0.6924 \quad (\text{triangle-free graphs})
    \]
  \item Matches known analytical results
  \end{itemize}
\end{frame}

%------------------------------------------------

\begin{frame}{Interpretation}
  \begin{itemize}
  \item $p=1$ captures:
    \begin{itemize}
    \item Short-range correlations
    \item Local structure of graph
    \end{itemize}
  \item Increasing $p$:
    \begin{itemize}
    \item Builds long-range correlations
    \item Approaches adiabatic limit
    \end{itemize}
  \end{itemize}
\end{frame}

%================================================
\section{QAOA as Trotterized Adiabatic Evolution}
%================================================

\begin{frame}{Adiabatic Evolution}
  \begin{itemize}
  \item Time-dependent Hamiltonian:
    \[
    \hat{H}(t) = (1-s(t))\hat{H}_M + s(t)\hat{H}_C
    \]
  \item Ground state evolves:
    \[
    |\psi(T)\rangle = \mathcal{T}
    \exp\left(-i\int_0^T \hat{H}(t) dt\right)|s\rangle
    \]
  \end{itemize}
\end{frame}

%------------------------------------------------

\begin{frame}{Trotterization}
  \begin{itemize}
  \item Discretize time:
    \[
    U(T) \approx \prod_{k=1}^p
    e^{-i\Delta t \hat{H}_C}
    e^{-i\Delta t \hat{H}_M}
    \]
  \item Compare with QAOA:
    \[
    \gamma_k \sim s(t_k)\Delta t, \quad
    \beta_k \sim (1-s(t_k))\Delta t
    \]
  \end{itemize}
\end{frame}

%------------------------------------------------

\begin{frame}{Variational Generalization}
  \begin{itemize}
  \item QAOA replaces schedule with free parameters:
    \[
    (\gamma_k,\beta_k) \text{ optimized}
    \]
  \item Advantages:
    \begin{itemize}
    \item No need for slow evolution
    \item Adaptation to hardware constraints
    \end{itemize}
  \end{itemize}
\end{frame}

%------------------------------------------------

\begin{frame}{Connection to Physics}
  \begin{itemize}
  \item QAOA $\leftrightarrow$ discrete quantum quench sequence
  \item Equivalent to:
    \begin{itemize}
    \item Digital adiabatic evolution
    \item Variational time evolution
    \end{itemize}
  \item Links to:
    \begin{itemize}
    \item Quantum control theory
    \item Many-body dynamics
    \end{itemize}
  \end{itemize}
\end{frame}

%------------------------------------------------

\begin{frame}{Limit $p \to \infty$}
  \begin{itemize}
  \item QAOA converges to adiabatic algorithm:
    \[
    \lim_{p \to \infty} |\psi_p\rangle \rightarrow |\psi_{\text{ground}}\rangle
    \]
  \item Trade-off:
    \begin{itemize}
    \item Large $p$ → accuracy
    \item Small $p$ → NISQ feasibility
    \end{itemize}
  \end{itemize}
  \end{frame}
